\section{Semigroup derivational space inside action chains}\label{sec:semigroup}

\subsection{Clean action chains}

Let
\[
P=\langle G\mid R\rangle
\]
be a finite semigroup presentation with nonempty finite generator set $G$ and defining relations between nonempty words.  Use a one-sorted signature consisting of a base constant $a$, one constant $b_g$ for each generator $g\in G$, and a single binary operation $\alpha$.  Set $W_\varepsilon(x)=x$. For a nonempty word $w=g_1\cdots g_k$, define
\begin{equation}
W_w(x)=
\alpha(b_{g_1},
 \alpha(b_{g_2},\ldots,
  \alpha(b_{g_k},x)\ldots)).
\label{eq:wordchain}
\end{equation}
Let $\mathcal A(P)$ consist of the universal identities
\begin{equation}
W_u(x)=W_v(x)
\qquad((u,v)\in R).
\label{eq:actionidentities}
\end{equation}

\begin{theorem}[Clean-chain correspondence]\label{thm:cleanchain}
For nonempty words $u,v$,
\[
\boxed{u=_P v\iff
\mathcal A(P)\vdash W_u(a)=W_v(a).}
\]
Moreover, elementary applications of a defining semigroup relation correspond exactly to elementary equational rewrites between clean chains.
\end{theorem}

\begin{proof}
A factor replacement $puq\leftrightarrow pvq$ is obtained by substituting $x\mapsto W_q(a)$ into the identity $W_u(x)=W_v(x)$ and placing the result in the outer context $W_p[-]$.

Conversely, every $\alpha$-rooted subterm of a clean chain is a suffix chain.  Any instance of a defining identity that matches a clean chain therefore substitutes a clean suffix for $x$, and the surrounding context is a clean prefix.  Hence every elementary equational rewrite projects to one semigroup factor replacement. Unrestricted equational equality is generated by such rewrites, proving the reverse implication as well.
\end{proof}

For equal words define
\[
\spc_P(u,v)=
\min_{u=w_0\leftrightarrow\cdots\leftrightarrow w_k=v}
\max_i|w_i|,
\]
\[
s_P(n)=
\max_{\substack{u=_Pv\\|u|,|v|\le n}}
\spc_P(u,v).
\]
Trivial pairs are included, so $s_P(n)\ge n$ for $n\ge1$. The maxima exist because there are only finitely many endpoint words of bounded length. This is the derivational-space function of the finite semigroup presentation \cite{Olshanskii2013}.  On clean chains, syntactic depth is word length, so Theorem~\ref{thm:cleanchain} yields
\begin{equation}
\boxed{
\lambdarw_{\mathcal A(P)}(W_u(a),W_v(a))
=
\spc_P(u,v).}
\label{eq:chainspace}
\end{equation}

\begin{example}[Congruence reuse does not measure word space]\label{ex:chaingap}
Take $P=\langle p,x,y,z\mid x=zzzz,\ y=zzzz\rangle$. Then
\[
\lambda_{\mathcal A(P)}(W_{px}(a),W_{py}(a))=4,
\qquad
\spc_P(px,py)=\lambdarw_{\mathcal A(P)}(W_{px}(a),W_{py}(a))=5.
\]
For the first equality, derive $W_x(a)=W_y(a)$ through the depth-four witness $W_{zzzz}(a)$ and apply congruence with first argument $b_p$. Below depth four no nontrivial ground instance of a defining identity is available. For the second equality, the first rewrite from $px$ must produce $pzzzz$, of length five, and $px\leftrightarrow pzzzz\leftrightarrow py$ attains that bound. Thus the exact clean-chain correspondence with space cannot use $\lambda$ in place of $\lambdarw$.
\end{example}

\subsection{The typed variant}
In the natural two-sorted variant
\[
\alpha:\mathsf{Op}\times\mathsf{State}\to\mathsf{State},
\]
with generator constants of sort $\mathsf{Op}$ and one state constant, every closed state term is already a clean chain.  There is no skeleton padding and
\[
\deltarw(n)=s_P(n)
\]
exactly for $n\ge1$, while $\deltarw(0)=0$. Generator-sort terms are immutable constants; the isolated base state contributes only a reflexive query. This pointwise identity concerns sequential rewrites.

\subsection{Skeleton decomposition of arbitrary one-sorted terms}

Call a node of a ground term \emph{marked} when it is literally of the form $\alpha(b_g,q)$ for some generator $g$.  A \emph{chain block} is a maximal consecutive sequence of marked nodes obtained by repeatedly following the second child.  Each block carries a nonempty generator word.  Collapse every maximal word to an abstract block label while preserving its location, tail, every unmarked $\alpha$-node, and all constants.  The resulting tree is the \emph{skeleton}.

\begin{lemma}[Skeleton invariance]\label{lem:skeleton}
One elementary rewrite using a defining identity in $\mathcal A(P)$ preserves the skeleton, preserves all chain-block words except one, and changes that one word by exactly one semigroup factor replacement.  Conversely, every defining factor replacement inside a chain block is realized by one equational rewrite.
\end{lemma}

\begin{proof}
A match of $W_u(x)$ can start only at a marked node and follows marked second-child nodes through the nonempty word $u$.  The substitution for $x$ is the suffix following the matched factor.  Replacing $u$ by $v$ therefore changes only that factor of the same maximal block.

No rewrite changes an unmarked node into a marked one or conversely.  Constants are immutable.  A nonconstant first child remains nonconstant because both sides of every defining identity are $\alpha$-rooted, and a non-generator constant remains itself.  Hence the marked/unmarked status of every surrounding node is invariant.  The converse is the same matching observation read backwards.
\end{proof}

\begin{corollary}[Global equality decomposition]\label{cor:skeleton}
Two arbitrary ground action terms are equal modulo $\mathcal A(P)$ if and only if they have the same skeleton and the words on every pair of corresponding chain blocks are equal in $P$.
\end{corollary}

Thus the one-sorted ground word problem is a syntactic skeleton comparison followed by finitely many semigroup word-problem instances.

\subsection{The contextual-envelope theorem}

\begin{theorem}[Global rewrite-visibility envelope]\label{thm:envelope}
For every $n\ge1$,
\begin{equation}
\boxed{
\deltarw_{\mathcal A(P)}(n)
=
\max_{1\le m\le n}\bigl(n-m+s_P(m)\bigr)}
\label{eq:envelope}
\end{equation}
and hence
\begin{equation}
\boxed{
\deltarw_{\mathcal A(P)}(n)-n
=
\max_{1\le m\le n}\bigl(s_P(m)-m\bigr).}
\label{eq:overhead}
\end{equation}
\end{theorem}

\begin{proof}
For the lower bound, fix $m\le n$ and equal words $u,v$ realizing $s_P(m)$.  The base constant $a$ is syntactically distinct from all generator constants.  Define frozen contexts
\[
C_0[x]=x,
\qquad
C_{k+1}[x]=\alpha(a,C_k[x]).
\]
Every outer node is unmarked, so no defining identity matches at its root.  Any sequential rewrite chain between
\[
C_{n-m}[W_u(a)]
\quad\text{and}\quad
C_{n-m}[W_v(a)]
\]
therefore performs a semigroup derivation inside a frozen context.  Its maximum term depth is $n-m$ plus the maximum intermediate word length.  Hence the pair has rewrite visibility depth $n-m+s_P(m)$.  Maximizing over $m$ gives the lower bound.

For the upper bound, let equal ground terms $s,t$ have depth at most $n$.  By Corollary~\ref{cor:skeleton}, they have the same skeleton.  If there are no chain blocks, the terms are identical and depth $n$ suffices. Otherwise index the corresponding chain blocks by $i$, with source and target words $u_i,v_i$, and put
\[
m_i=\max\{|u_i|,|v_i|\}.
\]
For each block choose a semigroup derivation whose intermediate lengths are at most $s_P(m_i)$.

For a fixed skeleton, total term depth is a monotone function of the block lengths and is $1$-Lipschitz in each block length separately: increasing one block by $r$ can increase a root-to-leaf path by at most $r$.

First process all blocks with $|u_i|\ge|v_i|$.  Before and after each conversion the current block-length vector is coordinatewise no larger than the source vector, so the term depth is at most $n$; while block $i$ is being rewritten, its temporary increase beyond the baseline length $m_i=|u_i|$ is at most $s_P(m_i)-m_i$.  The total depth is therefore at most
\[
n-m_i+s_P(m_i).
\]
Then process all blocks with $|u_i|<|v_i|$.  At this stage the current vector, with unprocessed increasing blocks still at source length, is coordinatewise bounded by the target vector.  Comparing the active block with the baseline in which it already has target length $m_i=|v_i|$ gives the same bound.  Hence a derivation exists with maximum depth at most
\[
\max_i\bigl(n-m_i+s_P(m_i)\bigr)
\le
\max_{1\le m\le n}\bigl(n-m+s_P(m)\bigr).
\]
\end{proof}

Define the contextual envelope of any function $s(n)\ge n$ by
\begin{equation}
H(s)(n)=n+\max_{1\le m\le n}(s(m)-m).
\label{eq:H}
\end{equation}

\begin{proposition}[Minimal context-stable envelope]\label{prop:H}
$H$ is the least majorant of $s$ whose overhead $H(s)(n)-n$ is nondecreasing.  It is idempotent:
\[
H(H(s))=H(s).
\]
If $s_P(n)-n$ is already nondecreasing, then
\[
\deltarw_{\mathcal A(P)}(n)=s_P(n)
\]
pointwise.
\end{proposition}

\begin{proof}
The running maximum in \eqref{eq:H} is nondecreasing and dominates $s(n)-n$, so $H(s)\ge s$.  If $f\ge s$ and $f(n)-n$ is nondecreasing, then for every $m\le n$,
\[
f(n)-n\ge f(m)-m\ge s(m)-m,
\]
so $f(n)\ge H(s)(n)$.  Minimality gives idempotence.
\end{proof}

Since $s_P$ is nondecreasing and $s_P(n)\ge n$,
\begin{equation}
s_P(n)\le\deltarw_{\mathcal A(P)}(n)
\le n+s_P(n)\le2s_P(n).
\label{eq:spacebounds}
\end{equation}
Following the standard equivalence for space functions \cite{Olshanskii2013}, write $f\preceq g$ if some positive integer $c$ satisfies
\[
f(n)\le c\,g(cn)+cn
\]
for all $n$, and write $f\simeq g$ when both directions hold.  Then
\begin{equation}
\boxed{\deltarw_{\mathcal A(P)}\simeq s_P.}
\label{eq:spaceequivalence}
\end{equation}

\subsection{Complexity-spectrum consequences}

The exact envelope formula transfers realization results for semigroup space functions.

\begin{corollary}[Deterministic-space realization]\label{cor:realization}
Let $f(n)$ be the space-complexity function of a deterministic Turing machine.  There exists a finite one-sorted equational presentation using only constants, one binary operation $\alpha$, and finitely many universal action-chain identities such that
\[
\deltarw_E\simeq f.
\]
\end{corollary}

\begin{proof}
Olshanskii proves that such $f$ is equivalent to the space function of a finitely presented semigroup \cite[Corollary~1.4]{Olshanskii2013}.  Apply Theorem~\ref{thm:envelope} and \eqref{eq:spaceequivalence}.
\end{proof}

\begin{corollary}
There exists a finite action-chain presentation with PSPACE-complete ground word problem and polynomial rewrite-visibility function.
\end{corollary}

\begin{proof}
Olshanskii gives a finitely presented semigroup with PSPACE-complete word problem and polynomial space function \cite[Corollary~1.5]{Olshanskii2013}.  Clean chains give the hardness reduction.  Corollary~\ref{cor:skeleton} reduces arbitrary one-sorted action-term equality to a skeleton comparison and finitely many semigroup word-problem instances, solvable sequentially in polynomial space.  Theorem~\ref{thm:envelope} gives polynomial rewrite visibility, and \eqref{eq:tworesources} also bounds congruence visibility from above.
\end{proof}

Olshanskii also determines, up to the same equivalence, which power functions occur as semigroup space functions \cite[Corollary~1.6]{Olshanskii2013}; the same realization criterion transfers to rewrite-visibility functions of finite action-chain presentations.

These realization statements concern $\deltarw$. They neither identify congruence visibility $\delta$ with semigroup space nor identify a presentation resource with the optimal machine-space complexity of its word problem.

\subsection{Immediate exactness under depth-preserving identities}

At the opposite extreme, call an identity $\ell=r$ \emph{depth-preserving} if every ground substitution $\sigma$ satisfies
\[
\depth(\ell\sigma)=\depth(r\sigma).
\]
If every identity of a finite presentation is depth-preserving, then every elementary rewrite preserves the depth of the whole term in every context.  Therefore no derivation starting in $T^{(d)}$ needs to leave it, and
\[
\boxed{\delta_E(d)=\deltarw_E(d)=d\qquad\text{for all }d.}
\]
This contrasts with action-chain relations of unequal word length, whose rewrite visibility can realize the semigroup-space behavior above.

